
\subsection{Why the extended physics moves so little here}
\label{subsec:whynull}

Three of this paper's results are nulls: station-resolved hydraulics change dispatch cost
by under one percent, transport delay by exactly nothing, and the piecewise linearisation
by a few tenths of a percent. A null invites two readings. Either these phenomena are
genuinely inert in district-heating dispatch, or the assumptions under which they were tested
closed the channels through which they would have acted. The second reading is the
sceptical one; it is legitimate, and it deserves an answer more specific than a caveat.

Four assumptions each close a specific channel. A \emph{prescribed heating curve} fixes the
supply and return temperatures, so the optimiser cannot trade thermal loss against pumping
energy, the trade in which hydraulics would acquire economic weight. \emph{Precomputed
heat-pump performance} makes the coefficient of performance independent of the network
state, so a heat pump cannot respond to a locally depressed supply temperature by shifting
its output in time or place. \emph{Fixed capacities} remove siting and sizing, so spatial
structure cannot express itself through where capacity is placed, only through where heat
is routed. And \emph{hourly resolution} discretises away every transient shorter than the
time step: With $\Delta t = 1$\,h, the trunk transit time of about $35$\,min gives $k_p = 0$: sub-hourly
delays are not represented on the hourly grid, and the delay step therefore shifts nothing on
this network. This is the reason the measured delay effect in
Table~\ref{tab:linearisation} lies below the solver tolerance.

The honest conclusion is therefore conditional, and stated that way throughout: within
the deterministic, hourly, fixed-capacity, fixed-temperature planning regime, extended
network physics adds little to a fixed schedule's forward-evaluated cost and reveals no
hydraulic violation. That regime is not a corner case. It is
where the overwhelming majority of published district-heating dispatch models operate, and
the practical guidance the nulls support (do not spend modelling effort on hydraulic
detail before loss representation) applies wherever such models are used.

What raises this above a plausible story is that one of the four channels has been opened
and measured the result. Freeing the supply temperature, the lever most likely to break the
hydraulic null, does break it: pipe velocity becomes the binding constraint at a 20\,K
reduction and pumping energy rises by a factor of four
(Section~\ref{subsec:tsup}). Hydraulics move from a negligible \emph{cost} to a binding
\emph{constraint} the moment the degree of freedom exists. That is direct evidence for the
mechanism rather than an appeal to it, and it sets the expectation for the three channels
left unopened: each should convert a null into a constraint rather than into a large cost
term, because the physics enters dispatch as a feasibility boundary and not as a driver of
the schedule.

One asymmetry is worth stating plainly because it determines how far the null results can be
trusted. Relaxing these assumptions would generally give the extended physics \emph{more} room
to matter, although not necessarily monotonically. The null results should therefore be
interpreted as conservative \emph{within the tested regime}; outside that regime, they are
indicative rather than guaranteed. The loss-dominance result is more robust: the decomposition
is an exact algebraic identity based on the four optimised costs, and the loss term dominates
empirically across all 135 synthetic networks, which span two orders of magnitude in loss
burden. Its magnitude remains conditional on the tested systems, but within that regime it
does not depend on any single one of the four assumptions above. This result, rather than the
null results, therefore motivates the title.
